\section{A published-GRR proof of rank nine for the nonsplit Gaussian $E_8$ datum}
\label{sec:e8-grr}

\paragraph{Scope and provenance.}
This section proves that a Hurwitz component for the Gaussian scalar
$E_8$ branch datum has a dominant Weil-sixfold period map.  The proof
uses an exact positive Hodge intersection, not a comparison of source
and target dimensions.  It is independent of the separate analytic
Riemann--Hilbert certificate.  The two exact numerical answers and the
proposal to use a rooted tree recurrence were leads in the unrefereed
rStation manuscript \cite{E8GRR:lead66}.  No theorem, character packet,
or executable from that manuscript is an input here.  The character
data are reconstructed from the CHEVIE author's table, the recurrence
is proved below, and the implementation is new.

An important source distinction is necessary.  Jarvis--Kimura's
\emph{A representation-valued relative Riemann--Hurwitz theorem and
the Hurwitz--Hodge bundle}, arXiv:0810.2488v3, supplies a useful
comparison formula, but no journal publication of that manuscript
has been verified.  It is not the published GRR premise of this proof.
The premise is Tseng's published formula \cite[\S7.3, especially
(7.3.6.1) and (7.3.7.2), pp.~60--64]{E8GRR:Tseng}.  Its normalization
is explained explicitly below.  Positivity is the published theorem
\cite[Lemma~6.17, with its proof and Lemma~6.15]{E8GRR:BBT}.
The references to these lemma numbers use the 48-page author version,
pp.~42--44.  No Severn, Markman, or 032 theorem is assumed.

\subsection*{D-E8-GRR-01. The fixed local system on every counted component}

Let $K=\mathbb Q(i)$.  In the simple-root basis let $H$ be the $E_8$
Cartan matrix: its diagonal entries are two and its entries on the
edges
\[
 (1,3),(3,4),(2,4),(4,5),(5,6),(6,7),(7,8)
\]
are minus one; other off-diagonal entries are zero.  Put
$s_j=1-e_jH_{j,*}$, $W=\langle s_1,\ldots,s_8\rangle$, and
$G=\langle W,iI\rangle\subset U(H,K)$.  Thus
\[
 |W|=696729600,\qquad |G|=1393459200,\qquad G=W\sqcup iW.
\]
Here $-I\in W$.  We use the natural representation $\rho$ on $K^8$.
It is absolutely irreducible and its rational simple group-algebra
factor is $M_8(K)$.

Write $R$ for the reflection class, $C_2$ for $i$ times the Weyl class
$2D_4(a_1)$, and $C_1$ for $i$ times the Weyl class $E_6$.  In the
CHEVIE class order these are respectively real class66, imaginary
class6, and imaginary class57.  The branch classes, with all markings
ordered, are
\[
                 (C_1,C_2,R,R,R,R,R,R,R,R,R,R).                 \tag{E8.1}
\]
The complete class words and sizes used here are retained in
\texttt{e8\_grr/character\_input.json}.  In particular, the word for
the real class57 is $s_1s_4s_2s_3s_6s_5$.  The orders and spectral
data required for (E8.1) are
\[
\begin{array}{c|r|r|l|r}
\text{class}&\text{size}&\text{order}&\text{eigenvalue exponents in }[0,1)
 &\text{age}\\\hline
R&120&2&0^7,(1/2)&1/2\\
C_2&15120&2&0^4,(1/2)^4&2\\
C_1&4838400&12&1/6,1/4,1/4,1/3,7/12,2/3,5/6,11/12&4
\end{array}                                                     \tag{E8.2}
\]
The exact matrix reconstruction described in D-E8-GRR-04 checks these
data.  Class57 is carried to class58 by multiplication by $-I$.

Let $\mathcal H$ be the stack of smooth principal $G$-covers of the
punctured genus-zero base, completed at the twelve marked points
with local classes (E8.1).  Disconnected total curves are allowed.
Equivalently its fiber over an ordered configuration is the groupoid
of product-one tuples in the specified classes, modulo simultaneous
conjugation by $G$.  Its mass is therefore the number of such tuples
divided by $|G|$.  Let $\overline{\mathcal H}$ be its proper
admissible-cover, or balanced twisted-cover, compactification.  We
use the normalization with this ordinary groupoid mass.  Its coarse
moduli map to $\overline M_{0,12}$ is finite, and it has pure
dimension nine.  It may have
several components, and the argument below initially sums all of them.

For a tuple in (E8.1), let $J\subset G$ be its generated subgroup.
The first inertia has no eigenvalue one, so $V^J=0$.  The moving-rank
sum is
\[
                 D=8+4+10=22.
\]
The parabolic-cohomology dimension is consequently $D-2\cdot8=6$
over $K$.  Chevalley--Weil gives the two Hodge dimensions three:
the age sum for $V$ is $4+2+10/2=11$, and the age sum for its dual
is also $22-11=11$, so each dimension is $11-8=3$.
These statements hold for every tuple, whether or not $J=G$.

The signed Hermitian determinant formula D-DET-01 applies since
$V^J=0$; it does not require that the image be all of $G$.
For a reflection the local positive determinant quotient is one.
For $C_2$ a column basis of $\operatorname{im}(1-g)$ has Gram
determinant16 and $\det(1-g)=16$, again giving one.  For $C_1$ the
moving space is all of $V$, with $\det H=1$ and $\det(1-g)=6$.
Thus its quotient is $1/6$, whose Gaussian norm class is three.
It follows that all the counted components have signed determinant
class three and signature $(3,3)$.

There is an actual abelian sixfold on each component after a finite
level base change.  The rational $M_8(K)$-isotypic summand of the
Jacobian of the total curve has the form $B^8$ up to isogeny; a
primitive self-adjoint projector gives $B$.  One can use the
rank-one projector $(1-s_j)/2$ inside this simple group-algebra
factor.  It is self-adjoint for the positive $G$-invariant root form,
and hence for the induced Rosati involution.  The center supplies
the $K$-action on $B$, and the restriction of the Jacobian
polarization is positive.  The cohomology of $B$ is the multiplicity
space just computed, of rational dimension12.  For a disconnected
curve its Jacobian means the product of the Jacobians of its connected
components; all these assertions continue to hold.  Thus counting
disconnected covers introduces no virtual abelian variety.

\subsection*{D-E8-GRR-02. The degree-one Hodge class and its boundary normalization}

For any finite-order element $g$ put
\[
 a(g)=\frac12\sum_{j=0}^{r-1}m_{g,j}\frac jr\left(1-\frac jr\right),
 \qquad r=\operatorname{ord}(g),                               \tag{E8.3}
\]
where $m_{g,j}$ is the multiplicity of $e^{2\pi i j/r}$ in $\rho$.
This is a class function, and $a(g^{-1})=a(g)$.  If
$\mathcal F_\rho=\operatorname{Hom}_G(V_\rho,\pi_*\omega)$, write
$\lambda_\rho=c_1(\det\mathcal F_\rho)$ on
$\overline{\mathcal H}$.  The rank of this bundle is three.

For a reduced boundary component $D_\alpha$, denote by $g_\alpha$
its node holonomy and by $r_\alpha$ its order.  Classes $\psi_j$
below are cotangent classes of the \emph{coarse base curve}.
Then
\[
 \boxed{\lambda_\rho=\sum_{j=1}^{12}a(C_j)\psi_j
          -\sum_\alpha a(g_\alpha)r_\alpha[D_\alpha].}          \tag{E8.4}
\]
When a boundary stratum has multiple branches, the sum is taken on
its normalizations with the usual pushforwards.  Formula (E8.4)
therefore has an intrinsic interpretation on the stack.

\begin{proof}
Apply Tseng's published formula (7.3.7.2) to $X=BG$ and the vector
bundle associated to $V_\rho$.  Here $TBG=0$, so the virtual class is
the actual fundamental class.  All positive-degree rational Chern
characters of a vector bundle on $BG$ vanish.  Put
\[
 A_2(g)=\sum_jm_{g,j}B_2(j/r),\qquad
 B_2(x)=x^2-x+\frac16.
\]
Thus $A_2(g)/2=8/12-a(g)$.  In degree one, the main term of Tseng's
formula is $8\kappa_1/12$, and the marking term is
$-\sum_j A_2(C_j)\psi_j/2$.  The node term in that formula is
$r^2 A_2(g)/4$ times the pushforward from the node gerbe with a
choice of branch.  The node gerbe has degree $1/r$ over the reduced
boundary, and the choice of branch has degree two.  Its contribution
is therefore
\[
                       \frac r2 A_2(g)[D].                    \tag{E8.5}
\]
This calculation includes the factor $r^2$ displayed in Tseng's
formula; dropping that factor would give a different intersection
problem.

Tseng uses the stack with sections to the marking gerbes.  By his
Remark2.4.2, p.~19, this is the fiber product of the marking gerbes
over the stack without those sections.  Its degree is
$\prod_jr_j^{-1}$.  The weighted class in \S2.5.1, p.~21, multiplies
by $\prod_jr_j$, and hence gives exactly the ordinary-cover mass
specified in D-E8-GRR-01.  Formula (E8.5) is unchanged on descending
to that normalization.  This also checks that no marking factor is
missing from the final answer.

The associated orbibundle has
$H^q(\mathcal C,V_\rho)=(H^q(C,\mathcal O_C)\otimes V_\rho)^G$.
Its $H^0$ vanishes by $V^J=0$.  Serre duality identifies its $H^1$
with $\mathcal F_\rho^*$, so the degree-one Chern character of its
derived pushforward is precisely $\lambda_\rho$.  Combining the
three terms just computed yields
\[
 \lambda_\rho=\frac8{12}
  \left(\kappa_1-\sum_j\psi_j+\sum_\alpha r_\alpha[D_\alpha]\right)
  +\sum_ja(C_j)\psi_j
  -\sum_\alpha a(g_\alpha)r_\alpha[D_\alpha].                 \tag{E8.6}
\]
The expression in parentheses is zero.  For example, apply the same
published GRR formula to the trivial representation: the derived
pushforward on a genus-zero base is the trivial line, its first
Chern class is zero, and $A_2=1/6$ gives exactly this relation.
Equivalently it is the pullback of the genus-zero Mumford relation,
since a node smoothing has the local form $t=s^{r_\alpha}$.
This proves (E8.4) from the published formula.
\end{proof}

For later use, the normal class of a reduced boundary is
\[
 [D_\alpha]|_{D_\alpha}
       =-\frac{\psi_++\psi_-}{r_\alpha}.                     \tag{E8.7}
\]
Indeed the coarse smoothing parameter is $t=s^{r_\alpha}$ and its
normal line is the tensor product of the two coarse tangent lines.
The normal line of $s=0$ is its $r_\alpha$-th root.  The local
equation also proves that the ramification-weighted class
$r_\alpha[D_\alpha]$ has the ordinary smooth-tuple mass over the
corresponding base boundary: a reduced boundary stratum has that
mass divided by $r_\alpha$.  At several nodes the factors multiply.
These facts can also be checked against Jarvis--Kimura's
Proposition3.3, equations(30)--(38), pp.~14--16, but that preprint is
not required for the proof.

\subsection*{D-E8-GRR-03. A complete rooted-tree recurrence}

We prove the finite recurrence used for the ninth power of (E8.4).
The argument applies to any finite group and any inversion-invariant
weight $a$.  Let $\mathcal A=Z(\mathbb Q[G])$, with class basis
$K_c=\sum_{g\in c}g$, ordinary central multiplication, and trace
\[
                  \tau(x)=\frac{[1]x}{|G|}.
\]
The trace pairing satisfies
\[
 \tau(K_cK_d)=\begin{cases}|c|/|G|,&d=c^{-1},\\0,&d\ne c^{-1}.
 \end{cases}                                                   \tag{E8.8}
\]
It is the gluing pairing for covers.  To see this directly, choose
trivializations of the torsor at one base point of each component
of a nodal genus-zero tree.  Each vertex imposes that the product
of its outgoing holonomies is one.  Identifying the two fibers at
an edge matches inverse holonomies.  Changes of trivialization
on the vertices can be eliminated along a spanning tree, leaving
one simultaneous conjugation by $G$.  Thus gluing yields the mass
of product-one tuples divided by $|G|$.  In the class basis the
inverse pairing is $|G|/|c|$; contracting an edge cancels one
vertex normalization and is exactly ordinary multiplication in
$\mathcal A$.  This proves the assertion without a connectedness
assumption on the torsor.

Write $\mathcal W$ for the diagonal linear operator
$\mathcal W(K_c)=a(c)K_c$.  Fix one marking as a root.  Let
$\mathbf c=(c_1,\ldots,c_q)$ count the other marked leaves by their
classes and put $k=\sum c_j$.  The following quantities are central
vectors whose coefficients are exponential generating coefficients
in the leaf counts, namely divided by $\prod_jc_j!$.

The vector $T[\mathbf c,a]$ sums rooted stable subtrees with a bare
root flag carrying an insertion $\psi_{\rm root}^{a}$; the root
itself has no marking exponential and no attached edge.  Every
other marking carries its exponential, and every internal edge
carries the boundary contribution derived below.  The possible
root exponent is $0\le a\le k-2$.

For a child attached to a parent flag carrying exponent $b$, let
$L[\mathbf c,b]$ include the additional factor $1/b!$ from the
parent vertex's psi integral.  A single marked leaf of class $c$
then has
\[
              L[\mathbf e_c,b]=\frac{a(c)^b}{(b!)^2}K_c.       \tag{E8.9}
\]
One factorial comes from its exponential and one from that psi
integral.  For a child subtree with $k\ge2$ the rule is
\[
 L[\mathbf c,b]=-
    \sum_{a=0}^{k-2}
       \frac{\mathcal W^{a+b+1}T[\mathbf c,a]}
            {(a+b+1)a!(b!)^2}.                               \tag{E8.10}
\]
Let $F_m[\mathbf c,b]$ be the coefficient of leaf count
$\mathbf c$ and total parent exponent $b$ in $L^m/m!$, with
multiplication in $\mathcal A$.  Then
\[
 T[\mathbf c,a]=
       \sum_{m=a+2}^{k}\frac{(m-2)!}{a!}
                    F_m[\mathbf c,m-2-a].                   \tag{E8.11}
\]
Finally, if the root marking has class $c_0$, the top intersection
on the $(k-2)$-dimensional Hurwitz stack is
\[
 (k-2)!\prod_jc_j!\;
 \tau\left(K_{c_0}\sum_{a=0}^{k-2}
                     \frac{a(c_0)^a}{a!}T[\mathbf c,a]\right).
                                                                  \tag{E8.12}
\]

\begin{proof}
Expand $\exp(t\lambda)$ using (E8.4).  Distinct compatible boundary
divisors give a stable tree; incompatible boundary divisors have
empty intersection.  Consider an edge whose reduced divisor is
$D$, whose inertia order is $r$, and whose weight is $a(c)$.
If this divisor occurs $\ell\ge1$ times in the expansion, its
coefficient, before the reduced-stratum mass is taken, is
\[
 \frac{(-t a(c)r)^\ell}{\ell!}
       \left(-\frac{\psi_++\psi_-}{r}\right)^{\ell-1}
 =-r\,\frac{t^\ell a(c)^\ell}{\ell!}
                 (\psi_++\psi_-)^{\ell-1}.                   \tag{E8.13}
\]
The reduced-stratum mass has the compensating factor $1/r$ by
(E8.7).  Orienting the edge toward the root chooses one of its
two branches, so it also removes the two-to-one cut-edge convention
in GRR.  The edge series in the central-algebra state sum is
therefore exactly
\[
       -\frac{\exp(t a(c)(\psi_++\psi_-))-1}
                   {\psi_++\psi_-}.                         \tag{E8.14}
\]
The coefficient of $\psi_+^a\psi_-^b$ is
\[
       -\frac{t^{a+b+1}a(c)^{a+b+1}}{(a+b+1)a!b!}.           \tag{E8.15}
\]
This proves (E8.10), including its sign and both endpoint
factorials; the extra $1/b!$ is the parent-vertex factorial.
Inversion of the child holonomy is immaterial because $a(c)$ is
inversion invariant.

For completeness, the genus-zero psi integral at a vertex of
valence $m+1$ is
\[
 \int_{\overline M_{0,m+1}}\psi_0^a\prod_{j=1}^m\psi_j^{b_j}
   =\frac{(m-2)!}{a!\prod_jb_j!},\qquad
                  a+\sum_jb_j=m-2.                           \tag{E8.16}
\]
This follows by induction from the three-point integral one and
the string equation.  In fact at least one exponent is zero;
forget that marking and apply the string equation.  The sum of
the resulting numerators is the sum of the remaining exponents,
which is $m-2$, giving (E8.16).  The string equation itself follows
from $\psi_j=f^*\psi_j+D_{j,\mathrm{forgotten}}$; the exceptional
sections are disjoint and their normal classes are the negative
cotangent classes.  Thus (E8.16) needs no assumption about the
Hurwitz cover.

A stable root vertex has $m\ge2$ children.  Its children already
include their factors $1/b_j!$.  Applying (E8.16), and summing
unordered children by the exponential construction $L^m/m!$,
gives (E8.11).  Restoring the exponential of the actual root
marking adds $a(c_0)^a/a!$, explaining the second root factorial
in (E8.12).  Restoring the labels of same-class leaves multiplies
by $\prod_jc_j!$.  Since we expanded an exponential, extracting
the top power rather than its exponential coefficient multiplies
by $(k-2)!$.  The last gluing normalization is the trace $\tau$.
This proves (E8.12).

There is no infinite series issue.  A term of $T[\mathbf c,a]$
has Hodge degree $k-2-a$, and a term of $L[\mathbf c,b]$ has
degree $k-1+b$.  For total dimension $d$, one can discard
$b>d-k+1$.  Every child in (E8.11) has fewer leaves than its
parent, so the recurrence is acyclic and terminates.  These
observations also prove the truncations in the supplied program.
\end{proof}

\subsection*{D-E8-GRR-04. Exact evaluation from the published character table}

For the classes in (E8.1), (E8.3) gives
\[
                     a(C_1)=\frac{17}{24},\quad
                     a(C_2)=\frac12,\quad a(R)=\frac18.       \tag{E8.17}
\]
The finite character table used in the evaluation is the complete
112-class table for $W(E_8)$ in the CHEVIE author's
\texttt{weyle8.g}.  Its primary URL and immutable downloaded hash are
\begin{quote}\small
\url{https://raw.githubusercontent.com/jmichel7/gap3-jm/master/pkg/chevie/tbl/weyle8.g}\\
SHA256 \texttt{7dff885669b7b86c45acb4b3b144ce673f}\allowbreak
\texttt{b0440736bd81b221d963694048e254}.
\end{quote}
The published package and the author's manual are cited in
\cite{E8GRR:Michel,E8GRR:CHEVIE}.  The complete class classification
is also Carter's published Table11, pp.~52--54
\cite{E8GRR:Carter}.  Completeness and the character status of the
CHEVIE table are standard finite-group inputs; orthogonality alone
is not being claimed as a proof that an arbitrary array consists
of genuine characters.

\texttt{build\_character\_input.py} specializes the published
polynomial Hecke-character data at $v=q=1$.  Omitted alternate rows
in that table are recovered by the documented sign twist, using
$(-1)^{\ell(w)}$ on each class word.  The retained file contains
the resulting integer character table, the degrees, the class
sizes and words, the central-translation permutation, and both
sets of Hodge weights.  It redistributes mathematical data, not
the entire third-party library.

\texttt{check\_character\_input.py} can recheck the retained
input without fetching the external library.  It reconstructs
all112 class matrices from the displayed Cartan matrix and words;
checks their exact orders and orthogonality; factors every exact
characteristic polynomial into cyclotomic polynomials; shifts
all exponents by $1/4$ for the $iW$ coset; and recomputes all224
weights by (E8.3).  It checks the trace of the natural character
on every reconstructed matrix, the complete character
orthogonality identity, the sum of class sizes and of squared
degrees, and multiplication by the central negative identity.
The natural character is row68.

Only rational arithmetic in the $W$-class algebra is needed.
Every subtree has at most one $C_2$ leaf and all its other leaves
are real reflections.  Thus it belongs to either $\mathbb Q[W]$
or $i\mathbb Q[W]$, and no multiplication inside a subtree
multiplies two imaginary factors.  For $x=\sum_cx_cK_c$ in the
$W$-class algebra its Fourier coordinates and their inverse are
\[
 \widehat x_r=\sum_cx_c\frac{|c|\chi_r(c)}{\chi_r(1)},\qquad
 x_c=\frac1{|W|}\sum_r\chi_r(1)\chi_r(c)\widehat x_r.          \tag{E8.18}
\]
All these Weyl characters are real.  The program checks that
these two matrices multiply to the identity, and multiplies
central vectors coordinatewise in Fourier coordinates.  Applying
$\mathcal W$ in (E8.10) is done in the class coordinates, using the
weight for the real or imaginary coset as appropriate.

At the final root multiplication there are exactly two imaginary
factors.  If $z=-I$, then $(iw)(iv)=zwv$.  Therefore the final
identity coefficient is the coefficient at the negative of the
root Weyl class, multiplied by that class size.  For root class57
this is class58.  The final normalization is $|G|=2|W|$, not
$|W|$.  This explains the central-product reduction used by the
program without treating the224-character table as an opaque
input.

Running \texttt{exact\_hodge\_volume.py} with ten reflections,
one $C_2$ leaf and root $C_1$ evaluates (E8.9)--(E8.12) exactly.
The result is
\[
 \boxed{\int_{\overline{\mathcal H}}\lambda_\rho^9
                       =\frac{3009825}{64}>0.}                \tag{E8.19}
\]
Independently of the Hodge weights, the ordinary tuple mass is
\[
 \boxed{\tau(K_{C_1}K_{C_2}K_R^{10})=90296156160000.}          \tag{E8.20}
\]
The retained output records exact rational arithmetic, the checked
transform identity, and the input hash.  The numerical values
from the unrefereed lead are compared only after the calculation;
they are not input to it.  The program also permits exchanging
$C_1$ and $C_2$ as the root, giving a differently rooted evaluation
of the same intersection.  The supplied verification record
distinguishes this comparison from an independent recurrence.

\subsection*{D-E8-GRR-05. Canonical extension and actual rank nine}

At least one component of $\mathcal H$ has a Weil-sixfold period
map of generic rank nine.  Consequently it is dominant onto the
unitary moduli component determined by its rational polarized
$K$-space, of signature $(3,3)$ and signed determinant three.

\begin{proof}
Let $\mathcal H_\beta$ be any irreducible component.  Pass to a
finite level cover, take a projective resolution of its
compactification, and resolve the boundary to obtain a smooth
projective variety $S$ with a normal crossings divisor $D$ and
a generically finite map to $\overline{\mathcal H}_\beta$.
After a further finite cover if necessary, local monodromy is
unipotent.  The admissible curve family is semistable; semistable
resolution after the base change does not change its Hodge
bundle.  Marking connected components of a disconnected family
requires only a further finite cover.

We check that the GRR extension used above pulls back to the
canonical Hodge extension required by the positivity theorem.
For a one-parameter semistable family of curves the extension of
$F^1$ is $\pi_*\omega$.  This published identification is recalled
in the proof of Theorem2.1 of \cite[p.~242]{E8GRR:CD}, with the
canonical extension characterized by nilpotent residues.  Locally
at a node $xy=t$, the relative logarithmic generator is
$dx/x=-dy/y$, which is precisely the dualizing generator.  Thus
semistable resolution and base change add no divisor to the
determinant line.  Apply this one-parameter identification to a
transverse disk at the generic point of every prime boundary
divisor of $S$, including its exceptional divisors.  The Hodge
bundles and canonical extensions are locally free; the
identification over $S\setminus D$ therefore extends over $S$
because it and its inverse have no codimension-one poles on the
normal variety $S$.  This proves the equality of extensions,
not just equality over the open family.  Constant rational
group-algebra projectors commute with this identification.
The direct-sum argument applies to disconnected curves as well.

Let $\mathbb V_{\rm full}$ be the rational full isotypic
variation for $\rho\oplus\overline\rho$.  It has rational rank96,
and its $F^1$ has rank48.  The fixed representation factors have
complex dimension eight, so its extended determinant line $L$
has
\[
 c_1(L)=8(\lambda_\rho+\lambda_{\overline\rho})
                         =16\lambda_\rho.                    \tag{E8.21}
\]
The last equality holds in rational first Chern classes on the
compactification: the weights (E8.3), and hence (E8.4), are
unchanged on dualizing the representation.  This is stronger
than an equality of ranks or of open Chern forms.

Apply Bakker--Brunebarbe--Tsimerman Lemma6.17 to
$\mathbb V_{\rm full}$ on $S\setminus D$.  Its hypotheses hold:
the variation is polarized and real, $S$ is smooth projective,
the boundary has normal crossings, and its local monodromies
are unipotent.  The lemma says that $L$ is nef, and that it is
big exactly when the Hodge differential
\[
 T_{S\setminus D}\longrightarrow
       \operatorname{Hom}(F^1,\mathbb V_{\rm full,\mathcal O}/F^1)
                                                                  \tag{E8.22}
\]
is injective as a map of sheaves.  In weight one this is the
period differential.  Sheaf injectivity is equivalent to
injectivity at the generic point.  On a smooth projective
ninefold a nef line bundle is big exactly when its ninth
self-intersection is positive, as also used in the published
proof of Lemma6.15.

It follows that the contribution
$\int_{\overline{\mathcal H}_\beta}\lambda_\rho^9$ of every
component is nonnegative: its pullback is $16^{-9}c_1(L)^9$,
and generically finite alterations multiply intersections by
their positive degree.  Formula (E8.19) is the sum of these
nonnegative contributions, so at least one is strictly
positive.  On that component (E8.22) is generically injective.
The full variation and its primitive abelian sixfold have the
same period differential kernel, since the former is a fixed
multiplicity of the latter.  The primitive period map therefore
has rank nine.

Finally, a polarized $K$-space of signature $(3,3)$ has a
unitary period domain of complex dimension nine.  The
algebraic family of sixfolds just constructed gives a map to
its moduli space after level structure.  A rank-nine point
gives an analytic open image, hence Zariski-dense image in
that moduli component.  This proves dominance.  The source
and target dimensions enter only at this final step, after
the full period rank has been established.
\end{proof}

\subsection*{D-E8-GRR-06. The dominant component consists of connected covers}

Using the Gaussian reflection-normal-closure reduction D-GAUSS-03
proved earlier in this work, the dominant component obtained in
D-E8-GRR-05 has full generated group $G$.  Thus its total curves
are connected.  This conclusion does not claim that every component
with branch datum (E8.1) is dominant or connected.

\begin{proof}
We first justify the generic endomorphism assertion needed here.
On the full unitary period domain of signature $(3,3)$, the only
rational endomorphisms which are Hodge for every point are the
given scalars $K$.  Here is a direct verification.  Identify the
real vector space with $\mathbb C^3\oplus\mathbb C^3$ and the
Hermitian form with $\operatorname{diag}(I_3,-I_3)$.  One of its
compatible complex structures is
$J_0=\operatorname{diag}(iI_3,-iI_3)$.  Write a real endomorphism
as a complex-linear matrix $A$ plus a complex-antilinear matrix
$B$ followed by conjugation.  Commutation with $J_0$ forces
$A$ to be block diagonal and $B$ to be block off-diagonal.
The tangent directions to compatible complex structures are,
up to a nonzero real factor,
\[
       \begin{pmatrix}0&-iZ\\ i\overline Z^{,t}&0\end{pmatrix},
                   \qquad Z\in\operatorname{Mat}_3(\mathbb C).
\]
Commutation with all these directions implies
$A_+Z=ZA_-$ for every $Z$, hence $A_+=A_-=aI_3$.
For the off-diagonal blocks of $B$ it implies
$B_{12}Z^t=ZB_{21}$ for every $Z$.  Taking $Z=I$ gives
$B_{12}=B_{21}=B_0$; taking diagonal $Z$ makes $B_0$ diagonal;
taking the off-diagonal matrix units then makes every diagonal
entry zero.  Thus $B=0$.  The real commutant is the complex
scalars, and its intersection with rational endomorphisms is
exactly $K$.

There are only countably many rational endomorphisms.  For each
one outside $K$ its Hodge locus in the domain is a proper closed
analytic subset by the calculation just made.  A nonempty open
subset cannot be a countable union of such subsets.  The
dominant period map in D-E8-GRR-05 therefore has a very general
fiber $B$ with $\operatorname{End}^0(B)=K$.  In particular $B$
is simple.  This argument uses the actual open period image and
does not assume a generic Mumford--Tate theorem.

Pass to a finite cover of the Hurwitz component to identify the
generated subgroup $J\subset G$ up to a fixed conjugation, and
decompose the coefficient representation $V|_J$ over $K$.
The semisimple algebra
\[
 \operatorname{End}_{KJ}(V)=\prod_\alpha
                         M_{m_\alpha}(D_\alpha)
\]
acts on its parabolic multiplicity Hodge structure $P$ (with
the opposite-algebra convention if the multiplicity is written
as a Hom).  Every summand projector is a Hodge projector.
Since the full rational Hodge endomorphism algebra of $P$ is
$K$, exactly one factor acts nontrivially.  A nonzero
homomorphism from a simple algebra to a field is injective,
so that factor must have $m_\alpha=1$ and $D_\alpha=K$.
Consequently there is a unique active absolutely irreducible
$K$-summand $\tau$ of $V|_J$, of dimension at most eight,
occurring once, and
\[
       P=H^1_{\rm par}(\mathbb P^1\setminus\{p_j\},\tau).
                                                                  \tag{E8.23}
\]
All the other summands have zero parabolic cohomology.
The decomposition can be made orthogonal for the positive
finite-group form.  Thus (E8.23) identifies the actual polarized
variation, not merely its Hodge numbers.

The character field of $\tau$ is $K$.  Indeed its matrices are
over $K$, so the only possible proper character subfield is
$\mathbb Q$.  If its character were rational, $\tau$ and its
conjugate would be isomorphic over $K$: an intertwiner is
obtained by solving the intertwining linear equations over
$K$, and any nonzero intertwiner between these absolutely
irreducible representations is invertible.  Composing this
intertwiner with scalar conjugation gives an invertible
$K$-antilinear, rational $J$-equivariant coefficient map.
It induces an invertible rational Hodge endomorphism of (E8.23)
which does not commute with $i$.  This contradicts
$\operatorname{End}_{\rm Hdg}(P)=K$.  This argument also excludes
a rational-character quaternionic representation; no implicit
Schur-index-one assertion over $\mathbb Q$ is being made.

Each of the ten reflection inertias in (E8.1) acts on $\tau$
either trivially or as an involutive reflection, because its
moving space on all of $V$ has dimension one.  At least one
acts nontrivially.  Otherwise the local system in (E8.23)
would have only two possible nontrivial monodromies, which
are inverse to each other.  The first has no fixed vectors
even on $\tau$, as it has none on $V$.  If $a=\dim_K\tau$,
the moving-rank sum would then be $2a$, and the genus-zero
parabolic dimension would be $2a-2a=0$, contrary to
$\dim_KP=6$.

The finite image of $J$ on $\tau$ is therefore absolutely
irreducible, has character field $K$, and contains reflections.
Its actual parabolic variation is the dominant nonsplit
sixfold variation (E8.23).  D-GAUSS-03 applies, including its
induced-representation cases.  It says that every irreducible
component of its reflection normal closure which can remain
in this situation is the natural $E_8$ representation.
Such a component has dimension eight.  Since
$\dim_K\tau\le8$, equality is forced, $\tau=V|_J$, and the
reflection normal closure in $J$ contains a copy of $W(E_8)$
in the natural representation.

All reflections in $G$ belong to its real subgroup $W$.
The reflection normal closure just found is consequently a
subgroup of the original $W$.  Its order is the full order
$696729600$ of $W(E_8)$, so it is $W$.  The tuple also
contains its specified element of $iW$, hence $J=G$.
The generated subgroup is constant up to conjugacy on a
Hurwitz component.  Its principal cover is connected exactly
when that subgroup is $G$.  This proves the assertion.
\end{proof}

\paragraph{Precise conclusion.}
There exists a dominant family of connected $G$-covers whose
actual primitive abelian factors have $K=\mathbb Q(i)$,
signature $(3,3)$, and signed determinant class three.
D-E8-GRR-01--05 establish dominance without a connectedness
premise; D-E8-GRR-06 proves connectedness using the separately
proved reflection-normal-closure barriers.  The positive total
intersection does not identify a particular explicit Nielsen
tuple as lying in the dominant component.  The construction of
algebraic detecting cycles is a separate assertion.
